\subsubsection{Tool Sequence Patterns in Successful vs Unsuccessful Repairs}
\label{sec:tool-sequences}

To understand how repair strategies differ between successful and unsuccessful attempts, we analyze tool sequence patterns using a sliding window approach. We extract consecutive tool invocations of length $k$ from each repair trajectory and count pattern frequencies. Using window size $k=3$, we identify the most common 3-step sequences executed by agents during repair attempts. We focus on the two lowest-performing agents, \qwencode and \geminicli, comparing patterns in successful repairs (where all tests pass) versus unsuccessful repairs (where tests fail).

Figure~\ref{fig:qwen_pass_fail_sequences} presents the comparison for \qwencode, which executed 3,557 successful 3-step sequences and 10,877 unsuccessful sequences. Figure~\ref{fig:gemini_pass_fail_sequences} presents the comparison for \geminicli, which executed 3,587 successful sequences and 9,033 unsuccessful sequences. Each visualization shows the top five most frequent patterns stratified by repair outcome, measured as percentages of all 3-step sequences within that outcome class.

\paragraph{Qwen Code: Structural Navigation Overhead}

\qwencode exhibits navigation overhead in both successful and unsuccessful repairs, indicating a structural rather than behavioral failure mode. The most frequent pattern in successful repairs is NAVIGATE $\rightarrow$ TEST $\rightarrow$ NAVIGATE (4.13\%), while unsuccessful repairs show the same pattern at 2.93\%. NAVIGATE operations appear in three of five top patterns for successful repairs and four of five for unsuccessful repairs. The presence of navigation-heavy patterns even in successful repairs suggests that Qwen's workflow requires frequent directory traversal regardless of outcome. The agent cannot achieve repairs without this overhead, indicating that the navigation requirement is intrinsic to its implementation rather than a symptom of difficulty during failed attempts.

The similar pattern distributions between successful and unsuccessful repairs indicate that Qwen does not adapt its strategy based on repair progress. Both outcomes show comparable frequencies of WRITE $\rightarrow$ NAVIGATE $\rightarrow$ BUILD (2.64\% Pass vs 2.86\% Fail) and READ $\rightarrow$ READ $\rightarrow$ READ (3.32\% Pass vs 2.45\% Fail). This consistency suggests that success or failure depends more on external factors (bug complexity, available context) than on the agent's adaptive behavior during repair. The navigation overhead consumes computational budget uniformly across all repair attempts, reducing available actions for productive activities.

\paragraph{Gemini CLI: Behavioral Shift Toward Over-Modification}

\geminicli demonstrates a dramatic behavioral shift between successful and unsuccessful repairs, revealing an adaptive failure mode. Successful repairs are dominated by WRITE $\rightarrow$ BUILD $\rightarrow$ TEST (9.09\%), representing the ideal iterative repair cycle: modify code, compile, and validate. This pattern is 1.96 times more frequent than the second most common pattern (READ $\rightarrow$ WRITE $\rightarrow$ BUILD at 4.74\%). The high concentration indicates that when Gemini succeeds, it consistently follows productive modification-compilation-testing workflows.

Unsuccessful repairs exhibit a fundamentally different pattern distribution. The most frequent pattern becomes WRITE $\rightarrow$ WRITE $\rightarrow$ WRITE (10.37\%), representing consecutive modifications without intermediate validation. This pattern accounts for only 4.63\% of successful repairs but increases 2.24-fold to 10.37\% in unsuccessful repairs, a difference of +5.75 percentage points. The shift indicates that when Gemini encounters difficulty, it responds by batching modifications rather than testing incremental changes. The second most frequent pattern in unsuccessful repairs remains WRITE $\rightarrow$ BUILD $\rightarrow$ TEST (8.48\%), but its relative frequency decreases compared to successful repairs (9.09\%), suggesting that failed attempts deviate from the productive cycle toward accumulated modifications.

Additional patterns reinforce this behavioral shift. WRITE $\rightarrow$ READ $\rightarrow$ WRITE appears at 6.06\% in unsuccessful repairs compared to lower frequencies in successful attempts, indicating that failed repairs involve modification, inspection, and further modification without compilation or testing. The increased prevalence of write-heavy patterns in unsuccessful repairs suggests that Gemini interprets test failures as signals to modify more aggressively rather than to validate changes incrementally.

\paragraph{Contrasting Failure Modes}

The comparison reveals two distinct relationships between tool sequences and repair outcomes. \qwencode exhibits a \textit{structural constraint}: navigation overhead persists regardless of success or failure, fragmenting workflows uniformly across all repair attempts. The agent's low repair accuracy stems from this consistent inefficiency rather than maladaptive behavior during difficult repairs. In contrast, \geminicli exhibits \textit{behavioral adaptation}: the agent modifies its strategy in response to difficulty, shifting from validated iterative cycles (WRITE $\rightarrow$ BUILD $\rightarrow$ TEST) to batched modifications (WRITE $\rightarrow$ WRITE $\rightarrow$ WRITE) when repairs prove challenging. This adaptation is counterproductive, as accumulated modifications without intermediate feedback increase the likelihood of compounding errors.

These findings suggest that improving low-performing agents requires different interventions. \qwencode would benefit from architectural changes to reduce navigation requirements, as its behavior is consistent but structurally inefficient. \geminicli would benefit from reinforcement strategies that discourage modification batching during difficult repairs, as its adaptive behavior exacerbates rather than resolves challenges. The absence of a universal pattern indicates that agent improvement strategies must account for individual behavioral characteristics rather than applying uniform optimizations.
